% Options for packages loaded elsewhere
\PassOptionsToPackage{unicode}{hyperref}
\PassOptionsToPackage{hyphens}{url}
\documentclass[
]{article}
\usepackage{xcolor}
\usepackage[margin=2.5cm]{geometry}
\usepackage{amsmath,amssymb}
\setcounter{secnumdepth}{-\maxdimen} % remove section numbering
\usepackage{iftex}
\ifPDFTeX
  \usepackage[T1]{fontenc}
  \usepackage[utf8]{inputenc}
  \usepackage{textcomp} % provide euro and other symbols
\else % if luatex or xetex
  \usepackage{unicode-math} % this also loads fontspec
  \defaultfontfeatures{Scale=MatchLowercase}
  \defaultfontfeatures[\rmfamily]{Ligatures=TeX,Scale=1}
\fi
\usepackage{lmodern}
\ifPDFTeX\else
  % xetex/luatex font selection
\fi
% Use upquote if available, for straight quotes in verbatim environments
\IfFileExists{upquote.sty}{\usepackage{upquote}}{}
\IfFileExists{microtype.sty}{% use microtype if available
  \usepackage[]{microtype}
  \UseMicrotypeSet[protrusion]{basicmath} % disable protrusion for tt fonts
}{}
\makeatletter
\@ifundefined{KOMAClassName}{% if non-KOMA class
  \IfFileExists{parskip.sty}{%
    \usepackage{parskip}
  }{% else
    \setlength{\parindent}{0pt}
    \setlength{\parskip}{6pt plus 2pt minus 1pt}}
}{% if KOMA class
  \KOMAoptions{parskip=half}}
\makeatother
\usepackage{longtable,booktabs,array}
\newcounter{none} % for unnumbered tables
\usepackage{calc} % for calculating minipage widths
% Correct order of tables after \paragraph or \subparagraph
\usepackage{etoolbox}
\makeatletter
\patchcmd\longtable{\par}{\if@noskipsec\mbox{}\fi\par}{}{}
\makeatother
% Allow footnotes in longtable head/foot
\IfFileExists{footnotehyper.sty}{\usepackage{footnotehyper}}{\usepackage{footnote}}
\makesavenoteenv{longtable}
\setlength{\emergencystretch}{3em} % prevent overfull lines
\providecommand{\tightlist}{%
  \setlength{\itemsep}{0pt}\setlength{\parskip}{0pt}}
\usepackage{bookmark}
\IfFileExists{xurl.sty}{\usepackage{xurl}}{} % add URL line breaks if available
\urlstyle{same}
\hypersetup{
  pdftitle={SST 409: Stochastic Processes},
  hidelinks,
  pdfcreator={LaTeX via pandoc}}

\title{SST 409: Stochastic Processes}
\usepackage{etoolbox}
\makeatletter
\providecommand{\subtitle}[1]{% add subtitle to \maketitle
  \apptocmd{\@title}{\par {\large #1 \par}}{}{}
}
\makeatother
\subtitle{Lecture Notes on Random Variables}
\author{}
\date{}

\begin{document}
\maketitle

{
\setcounter{tocdepth}{2}
\tableofcontents
}
\emph{A complete transcription of the 15 handwritten notebook pages. All
equations, symbols and notation have been re-typeset in LaTeX.}

\emph{Abbreviations used in the original: r.v. = random variable; st.
pr. = stochastic process; g.f. = generating function; p.g.f. =
probability generating function; dist. = distribution; mtd. = method;
qty = quantity; prob. = probability.}

\emph{--- p.~1 ---}

\section{1. Stochastic Processes: Definition and
Examples}\label{stochastic-processes-definition-and-examples}

Here, we mean a set of r.v. \(\{X(t)\}\) depending on a real parameter
\(t\) which varies in a certain set of natural nos.

A st. pr. is denoted by

\[\{X(t),\ t \in T\}.\]

In many problems of theory of st. pr., the parameter \(t\) means time.
Though it can denote any other physical qty such as space in which the
process occurs.

\subsection{Exp. of st. pr.}\label{exp.-of-st.-pr.}

\textbf{1.} Random walk where we are concerned with the motion of a
particle in a straight line which moves in discrete jumps with certain
probabilities i.e.~at each stage, it moves one unit to the left or right
with probability \(p\) or \(1-p\). If initially the particle is at the
origin then its position \(X(n)\) after \(n\) jumps is a discrete r.v.
depending on the discrete parameter \(n\). Thus \(\{X(n),\ n \in T\}\)
is a st. pr.

\textbf{2.} The no. of telephone calls \(X(t)\) during the arbitrary
period of time \(t\) is a discrete random variable depending on the
continuous parameter \(t\). The family of r.v.'s \(\{X(t),\ t \in T\}\)
constitutes a st. pr. \(\left[T = (0, \infty)\right]\).

\emph{Read 3 \& 4.}

\emph{--- p.~2 ---}

\section{2. Generating Functions}\label{generating-functions}

\textbf{Def:} Let \(a_0, a_1, a_2, \ldots\) be a sequence of real
numbers. Using a variable \(s\) we may define a function

\[A(s) = a_0 + a_1 s + a_2 s^2 + \cdots = \sum_{k=0}^{\infty} a_k s^k \qquad \cdots (1)\]

If \(A(s)\) (power series) converges in some interval \(s_0 < s < s_1\),
then \(A(s)\) is called the g.f. of the sequence
\(a_0, a_1, a_2, \ldots\). The variable \(s\) itself has no particular
significance.

Now:

\[\frac{dA(s)}{ds} = a_1 + 2a_2 s + 3a_3 s^2 + 4a_4 s^3 + \cdots + k a_k s^{k-1} + \cdots\]

\[\left. \frac{dA(s)}{ds} \right|_{s=0} = a_1\]

\[\frac{d^2 A(s)}{ds^2} = 2a_2 + 3 \times 2\, a_3 s + 4 \times 3\, a_4 s^2 + \cdots + k(k-1) a_k s^{k-2} + \cdots\]

\[\left. \frac{d^2 A(s)}{ds^2} \right|_{s=0} = 2a_2 \quad \Rightarrow \quad a_2 = \frac{1}{2} \left. \frac{d^2 A(s)}{ds^2} \right|_{s=0}\]

\[\left. \frac{d^3 A(s)}{ds^3} \right|_{s=0} = 3!\, a_3 \quad \Rightarrow \quad a_3 = \frac{1}{3!} \left. \frac{d^3 A(s)}{ds^3} \right|_{s=0}\]

\[\left. \frac{d^k A(s)}{ds^k} \right|_{s=0} = k!\, a_k \quad \Rightarrow \quad a_k = \frac{1}{k!} \left. \frac{d^k A(s)}{ds^k} \right|_{s=0}\]

\subsection{Exp 1}\label{exp-1}

If \(a_k = 1\) for all \(n \ge 1\), \(\{a_k\} = \{1, 1, 1, \ldots\}\),
we have \(a_0 = 1\), \(a_1 = 1\), \(a_2 = 1, \ldots\)

Then the g.f. of \(\{a_k\}\) is given by

\[A(s) = \sum_{k=0}^{\infty} a_k s^k = a_0 + a_1 s + a_2 s^2 + a_3 s^3 + \cdots = 1 + s + s^2 + s^3 + \cdots = \frac{1}{1-s}, \qquad |s| < 1\]

\emph{--- p.~3 ---}

\subsection{Exp 2}\label{exp-2}

Let \(\{a_k\} = \left\{ \frac{1}{k!} \right\}\).

Then \(a_0 = 1\), \(a_1 = 1\), \(a_2 = \frac{1}{2!}\),
\(a_3 = \frac{1}{3!}, \ldots\)

\(\therefore\) The g.f. of \(\{a_k\}\) is

\[A(s) = \sum_{k=0}^{\infty} a_k s^k = a_0 + a_1 s + a_2 s^2 + a_3 s^3 + \cdots = 1 + s + \frac{s^2}{2!} + \frac{s^3}{3!} + \cdots = e^s\]

\subsection{Exp 3}\label{exp-3}

Let \(\{a_k\} = \left\{ \binom{n}{k} \right\}\) for fixed \(n\).

\(\therefore\) The g.f. of \(\{a_k\}\) is

\[A(s) = \sum_{k=0}^{\infty} a_k s^k = \sum_{k=0}^{n} \binom{n}{k} s^k = (1+s)^n\]

\subsection{Exp 4}\label{exp-4}

Let \(a_0, a_1, a_2, \ldots, a_n\) be a sequence of nos satisfying the
recurrence relation

\[a_n + a_{n-1} - 16 a_{n-2} + 20 a_{n-3} = 0 \qquad \text{for } n \ge 3\]

subjected to the initial values \(a_0 = 0\), \(a_1 = 1\) \&
\(a_2 = -1\). Obtain the g.f. of the above sq. \(\{a_n\}\).

\subsubsection{Soln}\label{soln}

Let \(A(s) = \sum_{n=0}^{\infty} a_n s^n\) be the g.f. of \(\{a_n\}\).

\[\sum_{n=3}^{\infty} a_n s^n + \sum_{n=3}^{\infty} a_{n-1} s^n - 16 \sum_{n=3}^{\infty} a_{n-2} s^n + 20 \sum_{n=3}^{\infty} a_{n-3} s^n = 0\]

\[\sum_{n=3}^{\infty} a_n s^n + s \sum_{n=3}^{\infty} a_{n-1} s^{n-1} - 16 s^2 \sum_{n=3}^{\infty} a_{n-2} s^{n-2} + 20 s^3 \sum_{n=3}^{\infty} a_{n-3} s^{n-3} = 0\]

\[A(s) - a_0 - a_1 s - a_2 s^2 + s\left[A(s) - a_0 - a_1 s\right] - 16 s^2 \left[A(s) - a_0\right] + 20 s^3 A(s) = 0\]

\emph{--- p.~4 ---}

\[A(s) - 0 - s + s^2 + s\left[A(s) - 0 - s\right] - 16 s^2 \left[A(s) - 0\right] + 20 s^3 A(s) = 0\]

\[A(s) \left[1 + s - 16 s^2 + 20 s^3\right] = s\]

\[\therefore \quad A(s) = \frac{s}{1 + s - 16 s^2 + 20 s^3}\]

\section{3. Probability Generating Function
(P.G.F.)}\label{probability-generating-function-p.g.f.}

Suppose \(X\) is a r.v. which assumes non-negative integral values
\(0, 1, 2, \ldots\) and that \(\Pr\{X = k\} = p_k\),
\(k = 0, 1, 2, \ldots\) and \(\sum_k p_k = 1\), then the probability
generating function \(P(s)\) of the sequence of probabilities
\(\{p_k\}\) is given by:

\[P(s) = \sum_{k=0}^{\infty} p_k s^k = p_0 + p_1 s + p_2 s^2 + \cdots + p_k s^k + \cdots\]

\[= \sum_{k=0}^{\infty} \Pr\{X = k\}\, s^k = F(s^x)\]

where \(F(s^x)\) is the expectation of the function \(s^x\) (a r.v.) of
the r.v. \(x\).

The series \(P(s)\) converges for at least \(-1 \le s \le 1\); clearly
\(P(1) = 1\).

\textbf{NB:} It is possible to have a sequence of probabilities which
does NOT form a p.d.f.

e.g.~\(p_0 = 0\), \(p_1 = \frac{1}{4}\), \(p_2 = \frac{1}{2}\),
\(p_3 = \frac{1}{3}\), \(p_4 = \frac{3}{4}\)

\[A(s) = p_0 + p_1 s + p_2 s^2 + p_3 s^3 + p_4 s^4 = 0 + \frac{1}{4}s + \frac{1}{2}s^2 + \frac{1}{3}s^3 + \frac{3}{4}s^4 = \frac{1}{4}s + \frac{1}{2}s^2 + \frac{1}{3}s^3 + \frac{3}{4}s^4\]

\(A(s)\) is a g.f. of \(\{p_k\}\) but NOT a p.g.f. since:

\[p_1 + p_2 + p_3 + p_4 \ne 1\]

\subsection{3.1 Determination of mean \& variance using the
p.g.f.}\label{determination-of-mean-variance-using-the-p.g.f.}

\[P(s) = \sum_{k=0}^{\infty} p_k s^k\]

\emph{--- p.~5 ---}

\[\therefore \quad P'(s) = \sum_k k\, p_k s^{k-1}\]

\[P''(s) = \sum_k k(k-1)\, p_k s^{k-2}\]

The expectation \(E(x)\) is given by

\[E(x) = \sum_k k\, p_k = P'(1)\]

We have

\[E\{x(x-1)\} = P''(1) = \sum_k k(k-1) p_k\]

and

\[E(x^2) = E\{x(x-1)\} + Ex = P''(1) + P'(1)\]

Hence

\[\operatorname{Var}(x) = E(x^2) - \left[E(x)\right]^2 = P''(1) + P'(1) - \left[P'(1)\right]^2\]

\section{4. The p.g.f. \& Moments of Standard
Distributions}\label{the-p.g.f.-moments-of-standard-distributions}

\subsection{(i) Bernoulli dist.}\label{i-bernoulli-dist.}

Let \(x\) have a Bernoulli dist. with parameter \(p\)
i.e.~\(\Pr\{X = k\} = p_k = p^k q^{1-k}\), \(q = 1-p\), \(k = 0, 1\).

The p.g.f. of \(x\) is given by

\[P(s) = \sum_{k=0}^{\infty} p_k s^k = \sum_{k=0}^{1} p^k q^{1-k} s^k = q + ps\]

\[\therefore \quad P'(s) = p \qquad\qquad E(x) = P'(1) = p\]

\[P''(s) = 0 \quad \Rightarrow \quad P''(1) = 0\]

\[\operatorname{Var}(x) = P''(1) + P'(1) - \left[P'(1)\right]^2 = 0 + p - p^2 = p\left[1-p\right] = pq\]

\subsection{(ii) Binomial dist.}\label{ii-binomial-dist.}

Let \(x\) have a binomial dist. with parameters \(n\) \& \(p\) i.e.

\[\Pr[X = k] = p_k = \binom{n}{k} p^k q^{n-k}, \qquad k = 0, 1, 2, \ldots, n, \qquad q = 1-p\]

\emph{--- p.~6 ---}

\(\therefore\) The p.g.f. of \(x\) is given by:

\[P(s) = \sum_{k=0}^{\infty} p_k s^k = \sum_{k=0}^{n} \binom{n}{k} p^k q^{n-k} s^k = (q + ps)^n\]

\[\therefore \quad P'(s) = n(q+ps)^{n-1} p = np\,(q+ps)^{n-1}\]

\[\therefore \quad E(x) = P'(1) = np\,(q+p)^{n-1} = np \qquad \text{since } q + p = 1\]

\[P''(s) = n(n-1) p^2 (q+ps)^{n-2}\]

\[P''(1) = n(n-1) p^2 (q+p)^{n-2} = n(n-1) p^2 \qquad \text{since } q + p = 1\]

\[\operatorname{Var}(x) = P''(1) + P'(1) - \left[P'(1)\right]^2\]

\[= n(n-1)p^2 + np - n^2 p^2\]

\[= n^2 p^2 - np^2 + np - n^2 p^2 = np(1-p) = npq\]

\subsection{(iii) Geometric dist.}\label{iii-geometric-dist.}

\[\Pr\{X = k\} = p_k = q^k p, \qquad p + q = 1, \qquad k = 0, 1, 2, \ldots\]

The p.g.f. of \(x\) is given by:

\[P(s) = \sum_{k=0}^{\infty} p_k s^k = \sum_{k=0}^{\infty} q^k p\, s^k = p \sum_{k=0}^{\infty} (qs)^k\]

\[= p\left(1 + qs + (qs)^2 + (qs)^3 + \cdots\right)\]

\[= p \left[ \frac{1}{1-qs} \right] = \frac{p}{1-qs}\]

\[P'(s) = \frac{pq}{(1-qs)^2} \qquad \therefore \quad E(x) = P'(1) = \frac{pq}{(1-q)^2} = \frac{pq}{p^2} = \frac{q}{p}\]

\[P''(s) = \frac{2pq^2}{(1-qs)^3} \qquad\qquad P''(1) = \frac{2pq^2}{(1-q)^3} = \frac{2pq^2}{p^3} = \frac{2q^2}{p^2}\]

\[\therefore \quad \operatorname{Var}(x) = P''(1) + P'(1) + \left[P'(1)\right]^2\]

\[= \frac{2q^2}{p^2} + \frac{q}{p} - \frac{q^2}{p^2}\]

\[= \frac{q^2}{p^2} + \frac{q}{p} = \frac{q^2 + pq}{p^2} = \frac{q(q+p)}{p^2} = \frac{q}{p^2} \qquad \text{since } q + p = 1\]

\textbf{(iv) Poisson dist. --- Read.}

\emph{--- p.~7 ---}

\section{5. Remarks on the p.g.f.}\label{remarks-on-the-p.g.f.}

\textbf{i)} When \(s = e^t\), then

\[P\left(e^t\right) = E\left(e^{tx}\right) \qquad \text{which is the Moment generating fn. of } x.\]

\textbf{ii)}

\[\int_0^1 P(s)\, ds = \int_0^1 \left( \sum_{k=0}^{\infty} p_k s^k \right) ds\]

\[= \int_0^1 \left( p_0 + p_1 s + p_2 s^2 + p_3 s^3 + p_4 s^4 + \cdots + p_k s^k \right) ds\]

\[= \int_0^1 p_0\, ds + \int_0^1 p_1 s\, ds + \int_0^1 p_2 s^2\, ds + \int_0^1 p_3 s^3\, ds + \cdots\]

\[= p_0 + \frac{p_1}{2} + \frac{p_2}{3} + \frac{p_3}{4} + \cdots = \sum_{k=0}^{\infty} \frac{p_k}{k+1} = E\left( \frac{1}{X+1} \right)\]

\subsection{Exp}\label{exp}

Let \(x\) be a r.v. with p.g.f. \(P(s)\). Find the p.g.f. of
\(Y = mx + n\), where \(m, n\) are integers \& \(m \ne 0\).

\subsubsection{Soln}\label{soln-1}

Let \(P_Y(s)\) be the p.g.f. of \(Y\). Then;

\[P_Y(s) = E\left(s^Y\right) = E\left[s^{mx+n}\right]\]

\[= E\left(s^{mx} s^n\right) = s^n E\left(s^{mx}\right) = s^n P\left(s^m\right)\]

\subsection{Remarks}\label{remarks}

\textbf{i)} We can obtain the probabilities \(p_k\) from the p.g.f.
\(p_k\) can be found from \(P(s)\) i.e.

\[p_k = \frac{1}{k!} \left[ \left. \frac{d^k P(s)}{ds^k} \right|_{s=0} \right]\]

\textbf{ii)} \(p_k\) is also given by the co-efficient of \(s^k\) in the
expansion of \(P(s)\) as a power series in \(s\). e.g.

\emph{--- p.~8 ---}

\textbf{i)}

\[P(s) = e^{-\lambda(1-s)} = e^{-\lambda} e^{\lambda s}\]

\[= e^{-\lambda} \left[ 1 + \lambda s + \frac{(\lambda s)^2}{2!} + \frac{(\lambda s)^3}{3!} + \cdots + \frac{(\lambda s)^k}{k!} + \cdots \right]\]

\[\therefore \quad p_k = \text{co-efficient of } s^k = e^{-\lambda} \frac{\lambda^k}{k!} \qquad \text{which is a Poisson Distribution.}\]

\section{6. Sum of Random Variables}\label{sum-of-random-variables}

\subsection{a) Fixed numbers}\label{a-fixed-numbers}

Let \(X\) \& \(Y\) be 2 independent non-negative integral valued random
variables. The sum \(z = x + y\) is a r.v.

Let \(P_z(s)\) be the p.g.f. of \(Z\). Then;

\[P_z(s) = E\left(s^z\right) = E\left(s^{x+y}\right)\]

\[= E\left(s^x\right) E\left(s^y\right) \qquad \text{since } X \text{ \& } Y \text{ are independent}\]

\[= P_x(s)\, P_y(s)\]

{[} \(P_x(s)\): p.g.f. of \(X\); \(P_y(s)\): p.g.f. of \(Y\) {]}

This result says that the p.g.f. of the sum of two independent r.v.'s
\(X\) \& \(Y\) is the product of the p.g.f. of \(X\) \& \(Y\).

The result extends to the case when we have \(n\) independent random
variables \(x_1, x_2, x_3, \ldots, x_n\).

That is the p.g.f. of \(S_n = x_1 + x_2 + \cdots + x_n\) is the product
of the p.g.f.'s of the individual r.v.'s.

\[P_{S_n}(s) = E\left[s^{S_n}\right] = E\left[s^{x_1 + x_2 + \cdots + x_n}\right]\]

\[= E\left(s^{x_1}\right) E\left(s^{x_2}\right) \cdots E\left(s^{x_n}\right) \qquad \text{since } x_1, x_2, \ldots, x_n \text{ are independent}\]

\[= P_{x_1}(s)\, P_{x_2}(s) \cdots P_{x_n}(s)\]

\emph{--- p.~9 ---}

If \(x_1, x_2, \ldots, x_n\) are also identically distributed, each with
p.g.f. \(P(s)\), then the p.g.f. of \(S_n = x_1 + x_2 + \cdots + x_n\)
is

\[\left[P(s)\right]^n\]

\subsection{Exp}\label{exp-5}

Let \(x_1\) \& \(x_2\) be indep. poisson r.v.'s with parameters
\(\lambda_1\) \& \(\lambda_2\) resp. Find the distribution of
\(Z = X_1 + X_2\).

\subsubsection{Soln}\label{soln-2}

Let \(P_{x_1}(s)\) \& \(P_{x_2}(s)\) be the p.g.f.'s of \(X_1\) \&
\(X_2\) respectively. Then;

\[P_{x_1}(s) = \sum_{k=0}^{\infty} p_k s^k = \sum_{k=0}^{\infty} e^{-\lambda_1} \frac{\lambda_1^k}{k!} s^k\]

\[= e^{-\lambda_1} \sum_{k=0}^{\infty} \frac{(\lambda_1 s)^k}{k!}\]

\[= e^{-\lambda_1} e^{\lambda_1 s} = e^{-\lambda_1 (1-s)}\]

and similarly; \(P_{x_2}(s) = e^{-\lambda_2 (1-s)}\)

\(\therefore\) The p.g.f. of \(Z\) is

\[P_z(s) = P_{x_1}(s)\, P_{x_2}(s) = e^{-(\lambda_1 + \lambda_2)(1-s)}\]

This is the p.g.f. of a poisson r.v. with parameters \(\lambda_1\) \&
\(\lambda_2\). This shows the sum of the two indep. poisson r.v.'s is
also a poisson r.v.

This result extends to the case of \(n\) indep. poisson r.v.'s. That is
the sum \(S_n = x_1 + x_2 + \cdots + x_n\) of \(n\) indep. poisson
r.v.'s is a poisson r.v. with mean

\[\sum_{i=1}^{n} \lambda_i\]

where \(\lambda_i\) is the mean of \(x_i\).

\subsection{b) Random Number}\label{b-random-number}

\textbf{Theorem:} Let \(x_i\), \(i = 1, 2, \ldots\) be identically \&
independently dist. (i.i.d.) random variables with \(P(x_i = k) = p_k\)
\& p.g.f.

\[P(s) = \sum_{k=0}^{\infty} p_k s^k, \qquad i = 1, 2, 3, \ldots\]

Let \(S_N = x_1 + x_2 + \cdots + x_N\) where \(N\) is a r.v. independent
of the \(x_i\)'s. Let the distribution of \(N\) be given by
\(P[N = n] = g_n\) \& p.g.f.

\emph{--- p.~10 ---}

of \(N\) be

\[G(s) = \sum_{n=0}^{\infty} g_n s^n.\]

Then the p.g.f., \(H(s)\), of \(S_N\) is given by the compound function
\(G\left[P(s)\right]\) i.e.

\[H(s) = \sum_{j=0}^{\infty} P(S_N = j)\, s^j = G\left[P(s)\right]\]

\[H'(s) = G'\left[P(s)\right] P'(s)\]

\[\therefore \quad E(S_N) = H'(1) = G'\left[P(1)\right] \cdot P'(1)\]

\[= G'(1)\, P'(1) = E(N)\, E(x_i)\]

\subsection{Exp}\label{exp-6}

Suppose \(N\) has a poisson dist. with mean \(\lambda\) \& \(x_i\),
\(i = 1, 2, \ldots, N\) are i.i.d. with p.g.f. \(P(s)\). Find the p.g.f.
of \(S_N = x_1 + x_2 + \cdots + x_N\).

\subsubsection{Soln}\label{soln-3}

p.g.f. of \(S_N\) is

\[H(s) = G\left[P(s)\right]\]

But p.g.f. of a poisson r.v. with mean \(\lambda\) is given by;

\[G(s) = e^{-\lambda(1-s)}\]

\[\therefore \quad H(s) = e^{-\lambda\left(1 - P(s)\right)}\]

The mean is given by;

\[E(S_N) = E(x_i)\, E(N) = \lambda E(x)\]

\section{7. Tailed or Cumulative
Probabilities}\label{tailed-or-cumulative-probabilities}

\subsection{i) Upper tail}\label{i-upper-tail}

Let \(X\) have a p.d.f. \(\Pr\{X = k\} = p_k\), \(k = 0, 1, 2, \ldots\)
with p.g.f.

\[P(s) = \sum_{k=0}^{\infty} p_k s^k,\]

and

\[q_k = \Pr\{X > k\} = p_{k+1} + p_{k+2} + \cdots, \qquad k = 0, 1, 2, \ldots\]

with generating funct.

\[\phi(s) = \sum_{k=0}^{\infty} q_k s^k \qquad \left( \sum_k q_k = 1 \right)\]

\[\therefore \quad \phi(s) = \sum_{k=0}^{\infty} q_k s^k\]

\[= q_0 + q_1 s + q_2 s^2 + q_3 s^3 + \cdots\]

\[= (p_1 + p_2 + p_3 + \cdots) + (p_2 + p_3 + p_4 + \cdots) s + (p_3 + p_4 + p_5 + \cdots) s^2 + (p_4 + p_5 + p_6 + \cdots) s^3 + \cdots\]

\emph{--- p.~11 ---}

\[= p_1 + p_2(1+s) + p_3\left(1 + s + s^2\right) + p_4\left(1 + s + s^2 + s^3\right) + \cdots\]

\[= p_1 + p_2(1+s) + p_3 \frac{\left(1 - s^3\right)}{(1-s)} + p_4 \frac{\left(1 - s^4\right)}{(1-s)} + \cdots\]

\[= \frac{p_1(1-s) + p_2\left(1 - s^2\right) + p_3\left(1 - s^3\right) + p_4\left(1 - s^4\right) + \cdots}{(1-s)}\]

\[= \frac{(p_1 + p_2 + p_3 + p_4 + \cdots) - \left(p_1 s + p_2 s^2 + p_3 s^3 + p_4 s^4 + \cdots\right)}{(1-s)}\]

\[= \frac{(p_0 + p_1 + p_2 + p_3 + p_4 + \cdots) - \left(p_0 + p_1 s + p_2 s^2 + p_3 s^3 + p_4 s^4 + \cdots\right)}{(1-s)}\]

\[\therefore \quad \phi(s) = \frac{1 - P(s)}{1 - s}, \qquad -1 < s < 1\]

\subsubsection{Alt. mtd.}\label{alt.-mtd.}

\[q_k = p_{k+1} + p_{k+2} + p_{k+3} + \cdots\]

\[q_{k-1} = p_k + p_{k+1} + p_{k+2} + p_{k+3} + \cdots\]

\[p_k = q_{k-1} - q_k, \qquad k \ge 1\]

Defining \(P(s) = \sum_{k=0}^{\infty} p_k s^k\) \&
\(\phi(s) = \sum_{k=0}^{\infty} q_k s^k\)

\[\Rightarrow \quad \sum_{k=1}^{\infty} p_k s^k = \sum_{k=1}^{\infty} q_{k-1} s^k - \sum_{k=1}^{\infty} q_k s^k\]

\[P(s) - p_0 = s \sum_{k=1}^{\infty} q_{k-1} s^{k-1} - \left[\phi(s) - q_0\right]\]

\[P(s) - p_0 = s\, \phi(s) - \phi(s) + q_0 \qquad\qquad q_0 = p_1 + p_2 + p_3 + \cdots\]

\[P(s) - (p_0 + q_0) = \phi(s)(s - 1)\]

\[\phi(s)(1 - s) = (p_0 + q_0) - P(s)\]

\[\phi(s)(1 - s) = 1 - P(s) \qquad \text{since } p_0 + q_0 = 1\]

\[\therefore \quad \phi(s) = \frac{1 - P(s)}{1 - s}, \qquad -1 < s < 1\]

\emph{--- p.~12 ---}

\subsection{2) Lower tail}\label{lower-tail}

Let \(q_k = \Pr\{X \le k\}\), \(k = 0, 1, 2, \ldots\) \&
\(p_k = \Pr\{X = k\}\), \(k = 0, 1, 2, \ldots\)

Defining \(P(s) = \sum_{k=0}^{\infty} p_k s^k\) \&
\(\phi(s) = \sum_{k=0}^{\infty} q_k s^k\)

\[\therefore \quad q_k = p_k + p_{k-1} + p_{k-2} + \cdots + p_1 + p_0\]

\[q_{k-1} = p_{k-1} + p_{k-2} + \cdots + p_1 + p_0\]

\[\therefore \quad q_k - q_{k-1} = p_k, \qquad k \ge 1 \qquad \Rightarrow (1)\]

Multiplying both sides of the eqn (1) above by \(s^k\) \& summing over
\(k\) we obtain;

\[\sum_{k=1}^{\infty} q_k s^k - \sum_{k=1}^{\infty} q_{k-1} s^k = \sum_{k=1}^{\infty} p_k s^k\]

\[\therefore \quad \phi(s) - q_0 - s\, \phi(s) = P(s) - p_0\]

\[\phi(s)(1 - s) = P(s) - p_0 + q_0\]

\[\phi(s)(1 - s) = P(s) \qquad \text{since } p_0 = q_0\]

\[\therefore \quad \phi(s) = \frac{P(s)}{1 - s}, \qquad -1 < s < 1\]

\emph{--- p.~13 ---}

\section{8. Bivariate Probability Generating
Functions}\label{bivariate-probability-generating-functions}

Suppose \(X\) \& \(Y\) are integral valued r.v.'s with joint prob. dist.

\[\Pr\{X = j, Y = k\} = p_{jk}, \qquad j = 0, 1, 2, \ldots; \quad k = 0, 1, 2, \ldots\]

and;

\[\sum_j \sum_k p_{jk} = 1\]

Then the joint p.g.f. of \(X\) \& \(Y\) is given by;

\[P(s_1, s_2) = \sum_{j=0}^{\infty} \sum_{k=0}^{\infty} p_{jk}\, s_1^{\,j} s_2^{\,k} \qquad \text{assuming convergence}\]

for the values of \(s_1\) \& \(s_2\) i.e.~\(|s_1| < 1\), \(|s_2| < 1\).

\[P(s_1, s_2) = E\left(s_1^{\,X} s_2^{\,Y}\right) \quad \Rightarrow (1)\]

If \(s_1 = s_2 = s\), then \((1)\) becomes;

\[P(s, s) = E\left(s^X s^Y\right) = E\left(s^{X+Y}\right)\]

if \(X\) \& \(Y\) are mutually indep. r.v.'s then \(P(s_1, s_2)\)
becomes;

\[P(s_1, s_2) = E\left(s_1^{\,X}\right) E\left(s_2^{\,Y}\right)\]

If \(X\) \& \(Y\) are i.i.d., with a common p.g.f. \(G(s)\), then we
have

\[P(s_1, s_2) = E\left(s_1^{\,X}\right) E\left(s_2^{\,Y}\right) = G(s_1) \cdot G(s_2) = \left[G(s)\right]^2\]

\subsection{P.g.f.'s of Marginal
Distributions}\label{p.g.f.s-of-marginal-distributions}

Let \(X\) \& \(Y\) be 2 r.v.'s, then;

\[P(s_1, s_2) = E\left(s_1^{\,X} s_2^{\,Y}\right) \quad \rightarrow \text{joint p.g.f.}\]

The p.g.f. of \(X\) is given by;

\[P(s_1, 1) = E\left(s_1^{\,X}\right)\]

Similarly, p.g.f. of \(Y\) is given by;

\[P(1, s_2) = E\left(s_2^{\,Y}\right)\]

\emph{--- p.~14 ---}

\section{9. Calculations of Means, Variances \& Co-variances of
Bivariate Prob. Dist. Functions using the Joint
p.g.f.}\label{calculations-of-means-variances-co-variances-of-bivariate-prob.-dist.-functions-using-the-joint-p.g.f.}

\[P(s_1, s_2) \quad \rightarrow \quad \text{joint p.g.f. } (x \text{ \& } y)\]

\[P(s_1, 1) = P(s_1) \quad \text{---} \quad \text{p.g.f. of } x\]

\[P(1, s_2) = P(s_2) \quad \text{---} \quad \text{p.g.f. of } y\]

\[\therefore \quad E(x) = P'(s_1)\big|_{s_1 = 1} = \left. \frac{\partial P(s_1, 1)}{\partial s_1} \right|_{s_1 = 1}\]

\[E(Y) = P'(s_2)\big|_{s_2 = 1} = \left. \frac{\partial P(1, s_2)}{\partial s_2} \right|_{s_2 = 1}\]

The second moment is:

\[E\left[X(X-1)\right] = \left. \frac{\partial^2 P(s_1, 1)}{\partial s_1^{\,2}} \right|_{s_1 = 1}\]

and

\[E\left[Y(Y-1)\right] = \left. \frac{\partial^2 P(1, s_2)}{\partial s_2^{\,2}} \right|_{s_2 = 1}\]

\[\sigma_x^2 = E\left[X(X-1)\right] + E(x)^2 - \left[E(x)\right]^2\]

\[\sigma_y^2 = E\left[Y(Y-1)\right] + E(Y) - \left[E(Y)\right]^2\]

\[\operatorname{cov}(x, y) = \sigma_{xy} = E(xy) - E(x)E(y)\]

but;

\[E(xy) = \left. \frac{\partial^2 P(s_1, s_2)}{\partial s_1 \partial s_2} \right|_{s_1 = s_2 = 1}\]

\subsection{Exp}\label{exp-7}

Consider a series of Bernoulli trials with prob. of success \(p\).
Suppose that \(X\) denotes the no. of failures preceeding the 1st
success \& \(Y\) the no. of failures following 1st success \& preceeding
2nd success. Then the joint prob. distribution of \(X\) \& \(Y\) is
given by;

\[p_{jk} = \Pr\{X = j, Y = k\} = q^{j+k} p^2; \qquad j, k = 0, 1, 2; \qquad p + q = 1\]

\emph{--- p.~15 ---}

\subsubsection{Soln}\label{soln-4}

The joint p.g.f. is

\[P(s_1, s_2) = \sum_{j=0}^{\infty} \sum_{k=0}^{\infty} p_{jk}\, s_1^{\,j} s_2^{\,k}\]

\[= \sum_{j=0}^{\infty} \sum_{k=0}^{\infty} q^{j+k} p^2\, s_1^{\,j} s_2^{\,k}\]

\[= p^2 \sum_{j=0}^{\infty} (qs_1)^j \sum_{k=0}^{\infty} (qs_2)^k\]

\[= p^2 \left[ \left(1 + qs_1 + (qs_1)^2 + \cdots\right) \left(1 + qs_2 + (qs_2)^2 + \cdots\right) \right]\]

\[P(s_1, s_2) = p^2 \left( \frac{1}{1 - qs_1} \right) \left( \frac{1}{1 - qs_2} \right) = \frac{p^2}{(1 - qs_1)(1 - qs_2)}\]

\[P(s_1) = P(s_1, 1) = \frac{p^2}{(1 - qs_1)(1 - q)} = \frac{p}{1 - qs_1}, \qquad \text{since } q + p = 1\]

\[\therefore \quad P'(s_1) = \frac{pq}{(1 - qs_1)^2}\]

\[\therefore \quad E(x) = P'(1) = \frac{pq}{(1 - q)^2} = \frac{pq}{p^2} = \frac{q}{p}\]

\[P''(s_1) = \frac{2pq^2}{(1 - qs_1)^3}\]

\[P''(1) = \frac{2pq^2}{(1 - q)^3} = \frac{2pq^2}{p^3} = \frac{2q^2}{p^2}\]

\[\therefore \quad \operatorname{Var}(x) = P''(1) + P'(1) - \left[P'(1)\right]^2\]

\[= \frac{2q^2}{p^2} + \frac{q}{p} - \frac{q^2}{p^2} = \frac{q^2}{p^2} + \frac{q}{p}\]

\[= \frac{q^2 + pq}{p^2} = \frac{q(q+p)}{p^2} = \frac{q}{p^2} \qquad \text{since } q + p = 1\]

Similarly;

\[E(y) = \frac{q}{p} \qquad \& \qquad \operatorname{Var}(y) = \frac{q}{p^2}\]

\[\frac{\partial P(s_1, s_2)}{\partial s_1} = \frac{p^2 q}{(1 - qs_1)^2 (1 - qs_2)}\]

\[\therefore \quad \frac{\partial^2 P(s_1, s_2)}{\partial s_1 \partial s_2} = \frac{p^2 q^2}{(1 - qs_1)^2 (1 - qs_2)^2}\]

\section{Appendix A. Notes on the
Original}\label{appendix-a.-notes-on-the-original}

\emph{The following are minor notational slips in the handwritten
original. They have been transcribed above exactly as written; the
corrected forms are given here for study purposes.}

\begin{enumerate}
\def\labelenumi{\arabic{enumi}.}
\item
  \textbf{p.~6 (Geometric distribution, variance).} The original writes
  \(\operatorname{Var}(x) = P''(1) + P'(1) + \left[P'(1)\right]^2\); the
  sign of the last term should be minus:
  \(\operatorname{Var}(x) = P''(1) + P'(1) - \left[P'(1)\right]^2\). The
  subsequent line,
  \(= \frac{2q^2}{p^2} + \frac{q}{p} - \frac{q^2}{p^2}\), is handled
  correctly and the final answer \(\operatorname{Var}(x) = q/p^2\) is
  right.
\item
  \textbf{p.~14 (Bivariate variances).} The original writes
  \(\sigma_x^2 = E[X(X-1)] + E(x)^2 - \left[E(x)\right]^2\); the first
  moment term should be \(E(x)\),
  i.e.~\(\sigma_x^2 = E[X(X-1)] + E(x) - \left[E(x)\right]^2\). (Compare
  p.~5, where \(E(x^2) = E\{x(x-1)\} + Ex\) is used correctly.) The
  companion formula for \(\sigma_y^2\) is written correctly.
\item
  \textbf{p.~8, item (i).} This is the continuation of the Poisson
  p.g.f. started on p.~7, item (ii) of ``Remarks''; the notes read out
  \(p_k\) as the coefficient of \(s^k\) in \(P(s) = e^{-\lambda(1-s)}\),
  giving the Poisson probabilities \(p_k = e^{-\lambda}\lambda^k/k!\).
\item
  \textbf{p.~14, example.} The original gives \(j, k = 0, 1, 2, \ldots\)
  (the ``\(\ldots\)'' is implied by the geometric-type structure);
  strictly, \(X\) and \(Y\) range over all non-negative integers.
\end{enumerate}

\section{Appendix B. Page Index (source
images)}\label{appendix-b.-page-index-source-images}

{\def\LTcaptype{none} % do not increment counter
\begin{longtable}[]{@{}
  >{\raggedright\arraybackslash}p{(\linewidth - 4\tabcolsep) * \real{0.2069}}
  >{\raggedright\arraybackslash}p{(\linewidth - 4\tabcolsep) * \real{0.4828}}
  >{\raggedright\arraybackslash}p{(\linewidth - 4\tabcolsep) * \real{0.3103}}@{}}
\toprule\noalign{}
\begin{minipage}[b]{\linewidth}\raggedright
Page
\end{minipage} & \begin{minipage}[b]{\linewidth}\raggedright
Source image
\end{minipage} & \begin{minipage}[b]{\linewidth}\raggedright
Content
\end{minipage} \\
\midrule\noalign{}
\endhead
\bottomrule\noalign{}
\endlastfoot
1 & \texttt{bg1.png} & Definition of a stochastic process; two examples
(random walk, telephone calls) \\
2 & \texttt{bg2.png} & Generating functions: definition, coefficient
extraction by differentiation, Exp 1 \\
3 & \texttt{bg3.png} & Exp 2 (\(1/k!\)), Exp 3 (\(\binom{n}{k}\)), Exp 4
(recurrence relation, set-up) \\
4 & \texttt{bg4.png} & Completion of Exp 4; probability generating
function; NB; mean \& variance set-up \\
5 & \texttt{bg5.png} & Mean/variance from the p.g.f.; Bernoulli and
Binomial distributions \\
6 & \texttt{bg6.png} & Binomial (contd.), Geometric distribution;
Poisson (to be read) \\
7 & \texttt{bg7.png} & Remarks: mgf, integral of \(P(s)\), p.g.f. of
\(Y = mx+n\), recovering \(p_k\) \\
8 & \texttt{bg8.png} & Poisson p.g.f. expanded; sum of independent
random variables (fixed number) \\
9 & \texttt{bg9.png} & i.i.d. case; sum of two Poisson variables;
statement of the random-sum theorem \\
10 & \texttt{bga.png} & Compound (random-sum) p.g.f. \(G[P(s)]\); mean;
Poisson \(N\); upper tail \\
11 & \texttt{bgb.png} & Upper tail continued (two methods) \\
12 & \texttt{bgc.png} & Lower tail \\
13 & \texttt{bgd.png} & Bivariate p.g.f.; independence; marginal
p.g.f.'s \\
14 & \texttt{bge.png} & Means, variances, covariance from the joint
p.g.f.; example set-up \\
15 & \texttt{bgf.png} & Example solved: joint p.g.f., marginals,
\(E(x)\), \(\operatorname{Var}(x)\), cross-derivative \\
\end{longtable}
}

\end{document}
